\section{LIT-CIALDINI: Robert Cialdini Research: Influence and Social Proof}
\label{app:lity}

This appendix integrates papers by Robert Cialdini on social influence,
persuasion, and compliance. Cialdini's work reveals six principles of influence that shape
human behavior in systematic and predictable ways: reciprocity, commitment, social proof,
authority, liking, and scarcity.

\subsection{Research Program Overview}

Robert Cialdini's research program addresses core behavioral economics themes:

\begin{itemize}
  \item \textbf{Social Proof}: How others' behavior influences our decisions
  \item \textbf{Reciprocity}: Obligation to return favors and treatment
  \item \textbf{Authority}: Compliance through perceived expertise
  \item \textbf{Consistency}: Commitment to previous decisions
  \item \textbf{Liking}: Preference for attractive, similar people
  \item \textbf{Scarcity}: Value increases when availability decreases
\end{itemize}

\subsection{Papers Integrated: 8 works}

\subsubsection{1984: Influence: The Psychology of Persuasion...}

\paragraph{Citation.}
Cialdini, Robert B. (1984). ``Influence: The Psychology of Persuasion.''
\textit{Quill}.

\paragraph{Core Finding.}
Six principles of influence: reciprocity, commitment, social proof, authority, liking, scarcity

\paragraph{9C Integration.}
\begin{itemize}[nosep]
  \item \textbf{Domain}: Social Influence
  \item \textbf{Dimension}: S
  \item \textbf{Context}: Social Context
  \item \textbf{Complementarity}: γ = 0.72
\end{itemize}

\subsubsection{1993: Influence: The Psychology of Persuasion (Revised)...}

\paragraph{Citation.}
Cialdini, Robert B. (1993). ``Influence: The Psychology of Persuasion (Revised).''
\textit{HarperBusiness}.

\paragraph{Core Finding.}
Compliance tactics work by triggering automatic mental shortcuts

\paragraph{9C Integration.}
\begin{itemize}[nosep]
  \item \textbf{Domain}: Social Influence
  \item \textbf{Dimension}: S
  \item \textbf{Context}: Social Context
  \item \textbf{Complementarity}: γ = 0.73
\end{itemize}

\subsubsection{2001: Influence: Science and Practice...}

\paragraph{Citation.}
Cialdini, Robert B. (2001). ``Influence: Science and Practice.''
\textit{Pearson}.

\paragraph{Core Finding.}
Consistency principle: people commit to choices and defend them

\paragraph{9C Integration.}
\begin{itemize}[nosep]
  \item \textbf{Domain}: Social Influence
  \item \textbf{Dimension}: S
  \item \textbf{Context}: Consistency Pressure
  \item \textbf{Complementarity}: γ = 0.68
\end{itemize}

\subsubsection{2001: A Universal Norm Against Lying Promotes Cooperatio...}

\paragraph{Citation.}
Cialdini, Robert B., Wosinska, Wilhelmina (2001). ``A Universal Norm Against Lying Promotes Cooperation.''
\textit{Journal of Economic Behavior & Organization}.

\paragraph{Core Finding.}
Universal honesty norm increases cooperation rates

\paragraph{9C Integration.}
\begin{itemize}[nosep]
  \item \textbf{Domain}: Cooperation
  \item \textbf{Dimension}: S
  \item \textbf{Context}: Norm Context
  \item \textbf{Complementarity}: γ = 0.64
\end{itemize}

\subsubsection{2006: Ethical Influence: Mastering the Psychology of Per...}

\paragraph{Citation.}
Cialdini, Robert B. (2006). ``Ethical Influence: Mastering the Psychology of Persuasion.''
\textit{Harvard Business Review}.

\paragraph{Core Finding.}
Ethical persuasion uses same principles but maintains integrity

\paragraph{9C Integration.}
\begin{itemize}[nosep]
  \item \textbf{Domain}: Persuasion
  \item \textbf{Dimension}: S
  \item \textbf{Context}: Ethical Context
  \item \textbf{Complementarity}: γ = 0.55
\end{itemize}

\subsubsection{2008: Yes!: 50 Scientifically Proven Ways to Be Persuasi...}

\paragraph{Citation.}
Goldstein, Noah J., Martin, Steve J., et al. (2008). ``Yes!: 50 Scientifically Proven Ways to Be Persuasive.''
\textit{Free Press}.

\paragraph{Core Finding.}
50 evidence-based persuasion techniques validated through research

\paragraph{9C Integration.}
\begin{itemize}[nosep]
  \item \textbf{Domain}: Persuasion
  \item \textbf{Dimension}: S
  \item \textbf{Context}: Social Proof
  \item \textbf{Complementarity}: γ = 0.69
\end{itemize}

\subsubsection{2009: Pre-suasion: When the First Word Can Be the Differ...}

\paragraph{Citation.}
Cialdini, Robert B. (2009). ``Pre-suasion: When the First Word Can Be the Difference.''
\textit{Harvard Business Review}.

\paragraph{Core Finding.}
Timing of message presentation affects persuasion effectiveness

\paragraph{9C Integration.}
\begin{itemize}[nosep]
  \item \textbf{Domain}: Attention
  \item \textbf{Dimension}: E
  \item \textbf{Context}: Attentional Context
  \item \textbf{Complementarity}: γ = 0.71
\end{itemize}

\subsubsection{2016: Pre-suasion: A Revolutionary Way to Influence and ...}

\paragraph{Citation.}
Cialdini, Robert B. (2016). ``Pre-suasion: A Revolutionary Way to Influence and Persuade.''
\textit{Simon & Schuster}.

\paragraph{Core Finding.}
Pre-persuasion: direct attention before message delivery

\paragraph{9C Integration.}
\begin{itemize}[nosep]
  \item \textbf{Domain}: Attention
  \item \textbf{Dimension}: E
  \item \textbf{Context}: Attention Allocation
  \item \textbf{Complementarity}: γ = 0.74
\end{itemize}

\subsection{Summary}

This appendix integrates 8 papers by Robert Cialdini
spanning research on social proof, reciprocity, authority.
The papers demonstrate systematic patterns in human behavior that have important
implications for policy, organization, and individual decision-making.

\subsection{References}

For complete reference details and citations, see the master bibliography
\texttt{bcm2\_master\_references.bib}.

\vspace{1em}
\hrule
\vspace{0.5em}

\noindent\textit{Cross-references:}
Appendix G (Central Glossary), Chapter 14 (Applications)
